\section{Introduction}

Historically, portfolio managers have relied on the sample covariance matrix to estimate the risks and diversification benefits associated with different assets. However, when the number of assets is large relative to the available return observations, the sample covariance matrix can contain substantial estimation error. Minimum-variance optimisation is particularly sensitive to this problem because it tends to exploit the most extreme covariance estimates, often producing highly concentrated and unstable portfolio weights. Consequently, portfolios that appear optimal in-sample may perform poorly out-of-sample.

Covariance-matrix shrinkage seeks to address this problem by combining the sample covariance matrix with a more stable, structured target. By shrinking extreme and unreliable estimates toward this target, the method can reduce estimation error and improve the robustness of portfolio allocation. Building on the work of Ledoit and Wolf, this essay explains the intuition and construction of shrinkage estimators, evaluates alternative approaches to reducing covariance-estimation error, and empirically investigates the performance of shrinkage-based portfolios using equities listed on the Johannesburg Stock Exchange. The empirical results show that linear shrinkage produced modest improvements in numerical stability, portfolio weights and out-of-sample performance, while nonlinear shrinkage remained close to the sample estimator.
